% Avsnitt 8.8, ur lectures/L08/appendix/b_exercises.md.

\section{Övningar}
\label{c8:sec:exercises}

\begin{aside}
\textbf{Så kontrollerar du ditt arbete.} Övning~\ref{c8:ex:binary} och \ref{c8:ex:moore101} är penna
och papper: de handlar om att få en tillståndsmaskin rätt \emph{innan} den blir kod, och inget i dem
kräver VHDL, GHDL eller kort.

Från Övning~\ref{c8:ex:mealy101} och framåt har varje övning som ber dig skriva en VHDL-modul en
utdelad självkontrollerande testbänk under \file{lectures/L08/exercises/}, och ingenting annat;
Övning~\ref{c8:ex:fastblink} ändrar en modul du redan har och kontrolleras för hand. Att sätta ihop
resten av katalogen är en del av arbetet: modulerna du skrev i tidigare kapitel
(\code{reset_sync}, \code{button_sync}, \code{sync}, \code{timer}) kopieras in från där du än skrev
dem. Varje övning skriver ut de exakta kommandon den behöver.

Övning~\ref{c8:ex:rx} och \ref{c8:ex:tx} är bokens capstones: de två ändarna av en och samma
seriella länk, en mottagare och en sändare. Var och en komponerar moduler från fyra olika kapitel.
Avsätt riktig tid för dem, och konstruera varje tillståndsdiagram på papper innan du skriver någon
VHDL. Gör mottagaren först: den är den svårare av de två, eftersom den måste återskapa den
bittajming den blir tilldelad snarare än att definiera den själv.
\end{aside}

% ------------------------------------------------------------------------------------------
\exerciseset{Tillståndsdiagram och tillståndstabeller}

\exercise{En maskin med binär tillståndsordning}{Krets}
\label{c8:ex:binary}
Konstruera en ny Mooremaskin med fyra tillstånd, $\{Q_1, Q_2\}$, som styr en lysdiod via en enda
synkroniserad, flankdetekterad puls \code{X} från en knapp (producerad precis som \code{X} i det
genomarbetade exemplet i \secref{c8:sec:byhand}).

Den skiljer sig från det genomarbetade exemplet på två sätt:
\begin{itemize}
  \item Tillståndsordningen är den vanliga binärräkningen
    $00 \to 01 \to 10 \to 11 \to 00$, \textbf{inte} den Graykod som används i
    \secref{c8:sec:byhand}.
  \item Lysdioden (\code{Y}) lyser i \textbf{två} tillstånd i stället för ett: närhelst $Q_1 = 1$
    (alltså i både \code{10} och \code{11}).
\end{itemize}

\xpart{a} Rita tillståndsdiagrammet: fyra tillstånd, en övergångspil per tillstånd för $X = 1$, och
en självloop per tillstånd för $X = 0$.

\xpart{b} Fyll i hela tillståndstabellen (alla åtta rader av $\{Q_1, Q_2, X\}$).

\xpart{c} Härled $Q_1^{+}$, $Q_2^{+}$ och \code{Y} med Karnaughdiagram, på samma sätt som
\secref{c8:sec:byhand} härleder dem för Graykodsversionen.

\xpart{d} Realisera grindnätet för hand: två D-vippor för $\{Q_1, Q_2\}$, grindar som beräknar
$Q_1^{+}$, $Q_2^{+}$ och \code{Y}, och bygg och simulera det i CircuitVerse, med
dubbelvippsynkronisering och flankdetektering av knappingången precis som i \secref{c8:sec:cv}.

\xpart{e} Verifiera den i CircuitVerse, mot din egen tillståndstabell. Det finns ingen testbänk för
den här maskinen, och det är inte meningen att det ska finnas: tillståndstabellen du härledde i
\textbf{b)} \emph{är} specifikationen, och att kontrollera en krets mot en specifikation du själv
arbetat fram är hela poängen med övningen.

Stega klockan för hand och kontrollera, i den här ordningen:
\begin{itemize}
  \item \textbf{Stegningsvägen.} Pulsa \code{X} en gång per steg och gå ett helt varv,
    $00 \to 01 \to 10 \to 11 \to 00$. Läs av det förväntade nästa tillståndet ur din tabell före
    varje flank. Läs av $\{Q_1, Q_2\}$ på vippornas utgångar efter flanken. Alla fyra övergångarna
    måste stämma.
  \item \textbf{Hållvägen.} Stega klockan med $X = 0$ i vart och ett av de fyra tillstånden och
    bekräfta att maskinen står kvar. Det är fyra kontroller till, och det är den halva folk hoppar
    över; en maskin som stegar rätt men inte håller är en maskin vars $X'$-termer är fel.
  \item \textbf{Utgången.} Kontrollera \code{Y} mot din tabell i varje tillstånd: hög i \code{10}
    och \code{11}, låg i \code{00} och \code{01}. Gör det som en egen genomgång, inte medan du
    tittar på tillståndsbitarna.
\end{itemize}

Var de skiljer sig åt säger dig vilken grind du ska titta på. Ett felaktigt $Q_1^{+}$ syns bara på
de övergångar där $Q_1$ skulle ha ändrats. Ett felaktigt \code{Y} syns med tillståndsföljden
fortfarande helt korrekt, vilket är varför det är värt en egen genomgång.

\note[Tips]Eftersom tillståndsordningen ändrades från Graykod till binär, räkna med en övergång
($01 \to 10$) där \emph{två} tillståndsbitar vänder på en gång. Enligt \secref{c8:sec:byhand} är det
harmlöst för maskinen själv, eftersom ingenting samplar tillståndsregistret mellan flankerna, och
CircuitVerses ideala simulering visar dig ingenting åt något håll. Vad det däremot påverkar är
\code{Y}, som här avkodas kombinatoriskt ur tillståndsbitarna: på riktig hårdvara kan \code{Y}
glitcha kort medan de två bitarna är skeva. Lägg märke till vilken övergång det är, och vad du
skulle göra åt det om \code{Y} drev något som brydde sig.

% ------------------------------------------------------------------------------------------
\exerciseset{Sekvensdetektering}

\exercise{En Mooredetektor för \code{101}}{Resonemang}
\label{c8:ex:moore101}
Konstruera en \textbf{Moore}maskin som detekterar sekvensen \code{1, 0, 1} när den kommer en bit per
klockcykel på en ingång \code{din}, och som höjer sin utgång \code{y} så snart träffens tredje bit
setts. Icke-överlappande duger här: efter en fullständig träff, börja leta efter ett nytt
\code{1, 0, 1} från början i stället för att återanvända de avslutande bitarna.

\xpart{a} Rita tillståndsdiagrammet: ett tillstånd per hur stor del av mönstret som matchats
hittills, en pil ut ur varje tillstånd för \code{din = 0} \emph{och} för \code{din = 1} (en
sekvensdetektor är aldrig sysslolös, den antingen stegar framåt eller faller tillbaka), och markera
vilket tillstånd som driver \code{y = 1}.

\xpart{b} Fyll i tillståndstabellen, en rad per kombination av (tillstånd, \code{din}).

\xpart{c} Följ den. Välj en sexbitarssekvens för \code{din} som innehåller minst en träff, och
anteckna, cykel för cykel: tillståndet före flanken, värdet på \code{din}, och utgången \code{y}.

\xpart{d} Fällan i den här maskinen är pilarna tillbaka, inte de framåt. Säg för varje tillstånd vad
som händer på den ingång som \emph{inte} för träffen vidare: vilka av dem går tillbaka till starten,
och finns det någon ingång som borde skicka dig tillbaka till ett \emph{delvis} matchat tillstånd i
stället för hela vägen till början? Motivera ditt svar för det här mönstret.

\note[Tips]Räkna dina tillstånd innan du ritar dem. ”Inget matchat än”, ”sett \code{1}”, ”sett
\code{10}”, ”sett \code{101}” är den naturliga uppdelningen, och det sista av dem är det som driver
\code{y}.

\textbf{Spara ditt svar.} Övning~\ref{c8:ex:mealy101} implementerar samma detektor i VHDL, som en
Mealymaskin \textemdash{} som, enligt \secref{c8:sec:mealy}, behöver ett tillstånd färre än det du
just ritade. Att lista ut vilket av dina tillstånd som försvinner, och varför, är en del av den
övningen.

% ------------------------------------------------------------------------------------------
\exerciseset{Ändliga tillståndsmaskiner i VHDL}

\exercise{Ett fjärde tillstånd i \code{fsm_led}}{VHDL}
\label{c8:ex:fastblink}
Utöka \code{fsm_led} med ett fjärde tillstånd, \code{STATE_FAST_BLINK}, som blinkar med 20~ms i
stället för \code{STATE_BLINK}:s 100~ms. Följ det genomarbetade exemplet i \secref{c8:sec:vhdl}.

Arbeta på en \textbf{kopia} av modulen, inte på \file{lectures/L08/fsm_led/fsm_led.vhd} självt.
Kopiera först hela katalogen till ett eget ställe, tillsammans med de tre egna moduler den
instansierar:

\begin{shell}
cp -r lectures/L08/fsm_led /tmp/fsm_led_fast
cp lectures/L04/exercises/reset_sync/reset_sync.vhd /tmp/fsm_led_fast     # the three you wrote
cp lectures/L04/exercises/button_sync/button_sync.vhd /tmp/fsm_led_fast   # yourself
cp lectures/L07/exercises/timer/timer.vhd /tmp/fsm_led_fast
cd /tmp/fsm_led_fast
\end{shell}

Det finns ingen testbänk att köra här, men maskinen med fyra tillstånd ska ändå gå att analysera,
och det är värt att göra efter varje ändring: det fångar en saknad \code{case}-gren, eller ett
\code{state_t}-värde du lagt till i den ena processen men inte den andra.

\begin{shell}
ghdl -a --std=93 reset_sync.vhd button_sync.vhd timer.vhd fsm_led.vhd
\end{shell}

\begin{warning}
Skälet att arbeta på en kopia är värt att känna till, eftersom det är en verklig fallgrop och inte
sysselsättning för dess egen skull. \code{fsm_led} har en utdelad referenstestbänk,
\file{fsm_led_tb.vhd}, och Övning~\ref{c8:ex:readfsm} ber dig köra den. Den testbänken kontrollerar
cykeln med \textbf{tre} tillstånd: den trycker på knappen en gång från \code{STATE_BLINK} och
kontrollerar att lysdioden lyser fast, eftersom nästa tillstånd i maskinen som den är skriven är
\code{STATE_ON}. I din version med fyra tillstånd landar det trycket i \code{STATE_FAST_BLINK},
lysdioden blinkar fortfarande, och assertionen fallerar. Ingenting i felet talar om vilken av de två
övningarna det hör till. Håll referensmaskinen intakt så uppstår aldrig den frågan.
\end{warning}

Det finns ingen testbänk för maskinen med fyra tillstånd. Kontrollera den så som
\secref{c8:sec:byhand} lät dig kontrollera ett tillståndsdiagram: förutsäg tillståndsföljden för en
serie knapptryck, läs sedan tillbaka dina \code{case}-grenar och bekräfta att de stämmer.

Du behöver:
\begin{itemize}
  \item Lägga till det nya värdet i \code{state_t}.
  \item Lägga till en gren i varje \code{case}-sats, i både \code{STATE_PROCESS} och
    \code{LED_PROCESS}.
  \item Hantera de två blinktakterna, antingen genom att lägga till en andra timerinstans, eller
    genom att låta den befintliga timerns tickantal växla utifrån det aktuella tillståndet.
\end{itemize}

\exercise{\code{seq_detect_101_mealy}}{VHDL}
\label{c8:ex:mealy101}
Implementera detektorn för \code{1, 0, 1} från Övning~\ref{c8:ex:moore101} i VHDL, i en ny katalog
\file{seq_detect_101_mealy} \textemdash{} den här gången som en \textbf{Mealy}maskin.

Utgå från ditt tillståndsdiagram för Mooremaskinen och konvertera det:
\begin{itemize}
  \item Ditt sista tillstånd, ”sett \code{101}”, finns bara för att hålla utgången hög i en cykel.
  \item En Mealymaskin behöver det inte: utgångsvillkoret kan läsa \code{din} direkt i det tillstånd
    som redan matchat \code{1, 0}.
  \item Ta bort det tillståndet, och flytta in dess utgång i villkoret för \code{y}.
\end{itemize}

Följ sedan mönstret \code{type state_t is (...)} och \code{case} från \secref{c8:sec:vhdl} och
\ref{c8:sec:mealy}:
\begin{itemize}
  \item Återanvänd \code{reset_sync} för \code{reset_n}, precis som \file{seq_detect_mealy.vhd} gör.
    Det finns ingen knapp i den här designen, så \code{button_sync} instansieras helt enkelt inte.
  \item Se till att tilldelningen av utgången \code{y} är en enda konkurrent sats, utanför varje
    klockad process, som beror på både \code{state} och \code{din}, inte bara \code{state}. Annars
    har du i tysthet byggt en Mooremaskin i stället, och dess utgång kommer en cykel för sent
    \textemdash{} vilket är precis vad testbänken kontrollerar.
\end{itemize}

\modulefig[0.56]{lectures/L08/appendix/images/seq_detect_101_mealy.png}

\begin{selfcheck}
\textbf{Självkontroll.} Döp din entitet till \code{seq_detect_101_mealy}, med portarna \code{clock},
\code{reset_n}, \code{din} (in) och \code{y} (out), deklarerade i den ordningen för att matcha
\code{seq_detect_mealy}. En testbänk finns i \file{exercises/seq_detect_101_mealy/}; den
kontrollerar en icke-överlappande \code{101}-detektor. Din modul instansierar \code{reset_sync}, som
inte är utdelad där eftersom du skrev den i Övning~\ref{c4:ex:resetsync}: kopiera in din egen.

\begin{shell}
cd lectures/L08/exercises/seq_detect_101_mealy
cp ../../../L04/exercises/reset_sync/reset_sync.vhd .     # the one you wrote in L04
ghdl -a --std=93 reset_sync.vhd seq_detect_101_mealy.vhd seq_detect_101_mealy_tb.vhd
ghdl -e --std=93 seq_detect_101_mealy_tb
ghdl -r --std=93 seq_detect_101_mealy_tb --assert-level=error --stop-time=10ms
\end{shell}
\end{selfcheck}

% ------------------------------------------------------------------------------------------
\exerciseset{Att läsa den genomarbetade maskinen}

\exercise{Blinktakt, testbänk och tillståndsföljd}{Simulering}
\label{c8:ex:readfsm}
\code{fsm_led} demonstreras på DE0-CV (\secref{c8:sec:board}). Den här övningen kontrollerar att du
kan redogöra för det du såg, och den har en utdelad referenstestbänk du kan köra själv. Den
återanvänder dina två synkroniserare från \chapref{c4:ch} och din \code{timer} från
\chapref{c7:ch}, och ingen av dem är utdelad, så kopiera in dem först:

\begin{shell}
cd lectures/L08/fsm_led
cp ../../L04/exercises/reset_sync/reset_sync.vhd .     # the three you wrote yourself
cp ../../L04/exercises/button_sync/button_sync.vhd .
cp ../../L07/exercises/timer/timer.vhd .
ghdl -a --std=93 reset_sync.vhd button_sync.vhd timer.vhd \
                 fsm_led.vhd fsm_led_tb.vhd
ghdl -e --std=93 fsm_led_tb
ghdl -r --std=93 fsm_led_tb --assert-level=error --stop-time=10ms
\end{shell}

\xpart{a} Räkna ut blinktakten ur källkoden, inte ur demonstrationen. \code{fsm_led} instansierar
\code{timer} med \code{TIMER_TICK_COUNT} som har förvalet \code{5_000_000}, på en klocka på
\code{50 MHz}. Hur lång är en timeoutperiod? \code{STATE_BLINK} växlar lysdioden vid varje timeout.
Vilken \textbf{blink}frekvens ger det, och varför är den hälften av vad du kanske först skriver ned?

\xpart{b} Testbänken skriver över \code{TIMER_TICK_COUNT} med ett mycket mindre värde än förvalet.
Hitta den överskrivningen i \file{fsm_led_tb.vhd}. Förklara varför en testbänk inte rimligen kan
använda det riktiga värdet, och ungefär hur länge en simulering av en verklig blinkperiod skulle
behöva köra.

\xpart{c} Redogör för tillståndsmaskinens beteende:
\begin{itemize}
  \item Alla tre tillstånden går att nå åt båda hållen. Följ de tryck på \code{button_n(0)} och
    \code{button_n(1)} som besöker varje tillstånd och återvänder till \code{STATE_OFF}.
  \item Tryck på båda knapparna under samma klockcykel. Vad gör \code{fsm_led}, och vilka två rader
    i arkitekturen avgör det?
  \item \code{timer_enable} drivs ur \code{state}, så timern går bara i \code{STATE_BLINK}. Lägg
    märke till att \code{timer} \textbf{fryses} snarare än nollställs när den är avstängd
    (\secref{c7:sec:module}). Vad betyder det för ögonblicket då du återvänder till
    \code{STATE_BLINK} efter att ha lämnat det, och skulle det vara bättre eller sämre att hålla
    \code{timer_enable} hög permanent?
\end{itemize}

% ------------------------------------------------------------------------------------------
\exerciseset{Capstone: en seriell länk}

\exercise{\code{uart_rx8}: mottagaren}{VHDL}
\label{c8:ex:rx}
Det här är första halvan av bokens capstone, och den komponerar fyra kapitel till en enda design.
Övning~\ref{c8:ex:tx} bygger den andra änden av samma seriella länk, så stanna inte här. Ingenting i
den är en ny idé; allt handlar om att koppla ihop moduler du redan har, runt en tillståndsmaskin du
konstruerar själv:
\begin{itemize}
  \item \chapref{c4:ch}:s \code{reset_sync}, som synkroniserar reset.
  \item \chapref{c5:ch}:s \code{sync}, som synkroniserar ledningen \code{rx}, vilken kommer från
    någon annans kristall och är asynkron i precis den mening \chapref{c4:ch} avser.
  \item \chapref{c7:ch}:s \code{timer}, som genererar översamplingsticket.
  \item Det här kapitlets tillståndsmaskin, som ramar in alltihop och sköter sin egen skiftning och
    biträkning. Det byggde du som en modul i \chapref{c6:ch}; här är det fyra rader inuti maskinen
    som redan har de räknare den behöver, vilket är det vanligare sättet en mottagare skrivs på.
\end{itemize}

Skriv en entitet med namnet \code{uart_rx8}: en mottagare för en byte asynkron seriell data, i det
format en UART sänder.

\task{Ledningens format}
Ledningen \code{rx} bär ramar, och mellan ramarna ligger den i viloläge:
\begin{itemize}
  \item \textbf{Vila}: ledningen hålls hög.
  \item \textbf{Startbit}: ledningen går låg i exakt en bitperiod. Det är det som markerar början på
    en ram; det finns ingen separat klockledning, vilket är vad ”asynkron seriell” betyder.
  \item \textbf{Åtta databitar}, mest signifikant först, en bitperiod var.
  \item \textbf{Stoppbit}: ledningen går tillbaka hög i en bitperiod.
\end{itemize}

\lecfig[0.78]{lectures/L08/appendix/images/uart_frame.png}{En UART-ram: en ledning som vilar hög, en
startbit, åtta databitar med den mest signifikanta först, och en stoppbit, med mottagaren som
samplar mitt i varje bit.}{c8:fig:uartframe}

Ramen ovan bär \code{10110010}, som är den första byten testbänken skickar. Pilarna är de åtta
ögonblick då mottagaren tar upp en \textbf{databit}: en och en halv bitperiod efter den fallande
flanken, och sedan en gång per bitperiod därefter. Den tittar på ledningen även däremellan
\textemdash{} efter den fallande flanken, mitt i startbiten för att kontrollera att den fortfarande
är låg, och mitt i stoppbiten \textemdash{} men bara de åtta pilarna bidrar med bitar till den
mottagna byten.

Mottagaren måste återskapa bittajmingen enbart ur startbitens fallande flank. Det är hela problemet
den här övningen löser, och det är därför timern finns här.

\begin{aside}
\textbf{En medveten avvikelse från en riktig UART.} En riktig UART sänder sina databitar med den
\textbf{minst} signifikanta först, som \secref{c6:sec:sipo} påpekar. Den här övningen sänder dem med
den mest signifikanta först:
\begin{itemize}
  \item ramningen och tajmingen, som är vad de här två övningarna faktiskt handlar om, är identiska
    åt båda hållen.
  \item det är den riktning \secref{c6:sec:sipo}:s SIPO-idiom skiftar åt, så mottagarens skift blir
    en rad, och sändaren i Övning~\ref{c8:ex:tx} blir dess spegelbild.
  \item de utdelade testbänkarna sänder och förväntar sig MSB först för att matcha, så paret är
    konsekvent med sig självt.
  \item för att prata med en riktig UART skulle du vända skiftriktningen till
    \secref{c6:sec:shift}:s spegelvända form och inte ändra något annat. Värt att veta innan du
    kopplar det här till en dators serieport och undrar varför varje byte kommer bitvänd.
\end{itemize}
\end{aside}

\task{Översampling: sexton tick per bit}
Mottagaren måste hitta mitten av en bit vars början den bara kan sluta sig till. Tekniken heter
\textbf{översampling}: kör \code{timer} med \textbf{sexton gånger} baudhastigheten, så att det finns
sexton tick per bitperiod, och räkna dem.
\begin{itemize}
  \item Startbitens fallande flank nollställer tickräkningen till \code{0}. Timern går kontinuerligt,
    före, under och efter en ram; det är \emph{räkningen} som nollställs, aldrig timern.
  \item \textbf{Tick 8} ligger en halv bit senare, mitt i startbiten. Kontrollera att ledningen
    \emph{fortfarande är låg} där. Har den gått hög var det en glitch och ingen startbit, och
    mottagaren återgår till viloläge. Den omkontrollen är det som hindrar brus på en vilande ledning
    från att läsas som data.
  \item Därifrån är vart \textbf{sextonde} tick mitten på nästa bit: sampla ledningen och skifta in
    den. Åtta av dem är databitarna, och det nionde är stoppbiten.
\end{itemize}

Att sampla mitt i biten i stället för vid flankerna är det som ger mottagaren dess tajmingmarginal.
Sändarens klocka är en annan fysisk kristall än din, så de två driver isär under en ram; en
mottagare som samplade nära gränserna skulle läsa fel bit så snart de gjorde det. Sexton tick är det
konventionella valet, och sändaren i Övning~\ref{c8:ex:tx} går på samma tick.

\task{Entiteten}

\begin{booktable}{@{}p{11em}p{4.4em}p{3.4em}L@{}}
  \thead{Generic} & \thead{Typ} & \thead{Förval} & \thead{Beskrivning} \\
  \midrule
  \code{OVERSAMPLE_TICK_COUNT} & \code{natural} & \code{325} & Tickantalet för \code{timer} för
    \textbf{en sextondel} av en bitperiod. 9600 baud vid \code{50 MHz}; del \textbf{b)} ber dig
    härleda det. \\
\end{booktable}

\begin{booktable}{@{}p{5.6em}p{3.4em}p{10.8em}L@{}}
  \thead{Port} & \thead{Riktn.} & \thead{Typ} & \thead{Beskrivning} \\
  \midrule
  \code{clock} & in & \code{std_logic} & Systemklocka. \\
  \code{reset_n} & in & \code{std_logic} & Asynkron, aktiv låg; går genom \code{reset_sync}, som
    alltid. \\
  \code{rx} & in & \code{std_logic} & Den seriella ledningen, direkt från en pinne och därför
    \textbf{asynkron}. Synkronisera den till \code{rx_s2} innan något annat läser den. \\
  \code{data_out} & out & \code{std_logic_vector(7 downto 0)} & Den mottagna byten, giltig under den
    cykel då \code{byte_valid} är hög. \\
  \code{byte_valid} & out & \code{std_logic} & En puls på en cykel: en komplett, korrekt ramad byte
    ligger på \code{data_out}. \\
  \code{frame_err} & out & \code{std_logic} & En puls på en cykel: stoppbiten var inte hög, så ramen
    är trasig och ingen byte lämnas över. \\
\end{booktable}

\xpart{a} Konstruera tillståndsmaskinen \textbf{på papper först}, precis som \secref{c8:sec:byhand}
lät dig göra. Fyra tillstånd:
\begin{itemize}
  \item \code{STATE_IDLE}, som väntar på att ledningen ska gå låg.
  \item \code{STATE_START}, som vid tick 8 kontrollerar att ledningen fortfarande är låg.
  \item \code{STATE_DATA}, som samplar en bit vart sextonde tick, åtta gånger.
  \item \code{STATE_STOP}, som samplar stoppbiten sexton tick efter den sista databiten och väljer
    mellan \code{byte_valid} och \code{frame_err}.
\end{itemize}

Rita den innan du skriver någon VHDL, och på diagrammet:
\begin{itemize}
  \item Ge varje tillstånd sin \textbf{självloop}: pilen för ”ticket jag väntar på har inte kommit
    än, stå kvar”. Det är den pil som brukar saknas, och ett tillstånd utan den faller igenom dit
    \code{else} råkar peka.
  \item Skriv bredvid varje tillstånd \textbf{hur många tick det väntar på}, så att tickbokföringen
    ligger på papperet i stället för i huvudet.
  \item Markera var \code{ticks} \textbf{nollställs}. Timern själv stannar aldrig, så de
    nollställningarna är de enda ställen där mottagaren avgör vad ”början på en bit” betyder.
  \item Märk pilen ut ur \code{STATE_IDLE} med en \textbf{flank}, inte en nivå: \code{rx_s2} låg
    \emph{och} hög vid föregående tick. Det kräver en bit historik, så anteckna var du tänker hålla
    den.
\end{itemize}

\xpart{b} Räkna ut tajmingen innan du implementerar den:
\begin{itemize}
  \item DE0-CV:s klocka är \code{50 MHz} och 9600 baud betyder 9600 bitar per sekund. Hur många
    klockcykler är en bitperiod? Hur många är en sextondel av en?
  \item \code{timer} pulsar var \code{TICK_COUNT + 1} klockcykel (\secref{c7:sec:concept}). Vilket
    \code{OVERSAMPLE_TICK_COUNT} ger det dig? Det exakta svaret är inget heltal, så det finns två
    kandidater på var sin sida om det: räkna ut felet var och en av dem lämnar, och säg vilken av
    dem förvalet ovan är och varför.
  \item Det felet är en bråkdel av en procent, och det ackumuleras över en ram. Hur långt har din
    samplingspunkt tio bitperioder efter startflanken drivit från mitten av stoppbiten, mätt i
    översamplingstick? Hur många tick marginal hade den från början?
  \item Och nu skälet till att svaret är bekvämt: du omsynkroniserar vid varje startbit. Vad skulle
    behöva gälla för ramlängden för att det skulle sluta räcka?
\end{itemize}

\xpart{c} Implementera \code{uart_rx8}. Den instansierar tre moduler och gör resten själv:
\begin{itemize}
  \item \code{reset_sync} för reset. Allt annat i designen tar dess \code{reset_s2_n}.
  \item \textbf{\code{sync}}, modulen du skrev i Övning~\ref{c5:ex:sync}, på pinnen \code{rx}, som
    ger det \code{rx_s2} tillståndsmaskinen läser. Instansiera den med \code{SIZE = 1} och
    \code{PRESET = '1'}: en vilande ledning är hög, och en synkroniserare som kommer ur reset och
    läser \code{'0'} skulle för den här maskinen se ut som en startbit. \code{rx} kommer från en
    pinne som drivs av en sändare med sin egen kristall, så den är asynkron i precis den mening
    \chapref{c4:ch} avser, och varje regel därifrån gäller oförändrad. Lägg märke till att den tar
    en \code{std_logic_vector} med ett element på var sida, så du deklarerar ett par sådana och
    plockar ut \code{rx_s2} ur utgången.
  \item \code{timer}, med \code{OVERSAMPLE_TICK_COUNT} skickat rakt igenom, \code{enable} bundet
    till \code{'1'}, och dess \code{timeout} som din enda uppfattning om förfluten tid. Kalla den
    \code{baud_tick}. Timern frilöper: den stoppas aldrig, nollställs aldrig, och inget tillstånd rör
    den någonsin. Det maskinen fasar om i stället är sin egen räknare \code{ticks}.
  \item Ingen skiftregistermodul. Maskinen behöver ändå en biträknare, så den skiftar in ledningen
    direkt i en signal \code{frame} själv, en rad VHDL per samplad bit.
\end{itemize}

\task{Tillståndsmaskinen, tillstånd för tillstånd}
Två räknare: \code{ticks} som räknar \code{0} till \code{15} inom en bit, och \code{bit_count} som
räknar de åtta databitarna. Allt nedan händer \textbf{bara när \code{baud_tick = '1'}}; vid varje
annan klockflank står maskinen still. Eftersom timern frilöper är ett tick alltid på väg, i varje
tillstånd inklusive vila.

\begin{booktable}{@{}p{6.6em}p{7em}p{10em}L@{}}
  \thead{Tillstånd} & \thead{Väntar på} & \thead{Sedan} & \thead{Till} \\
  \midrule
  \code{STATE_IDLE} & fallande flank på \code{rx_s2} & nollställ \code{ticks} &
    \code{STATE_START} \\
  \code{STATE_START} & \code{ticks = 8} & \code{rx_s2} fortfarande låg? nollställ \code{ticks} &
    \code{STATE_DATA} \\
   & & \code{rx_s2} hög? det var en glitch & \code{STATE_IDLE} \\
  \code{STATE_DATA} & \code{ticks = 15} & skifta in \code{rx_s2} i \code{frame}, nollställ
    \code{ticks} & \code{STATE_DATA} tills 8 bitar, sedan \code{STATE_STOP} \\
  \code{STATE_STOP} & \code{ticks = 15} & \code{rx_s2} hög? \code{data_out <= frame}, pulsa
    \code{byte_valid} & \code{STATE_IDLE} \\
   & & \code{rx_s2} låg? pulsa \code{frame_err}, kasta byten & \code{STATE_IDLE} \\
\end{booktable}

Det finns inget ”klart”-tillstånd. Samplingen av stoppbiten är där ramen bedöms, så det är där byten
lämnas över, och maskinen går rakt tillbaka till viloläge.

\task{Varför vilotillståndet väntar på en flank i stället för en nivå}
Håll ett tick historik för \code{rx_s2}, och lämna vilotillståndet bara vid övergången från hög till
låg. Att i stället testa nivån, ”om \code{rx_s2} är låg, starta en ram”, ser likvärdigt ut och
klarar allt utom det enda fall som spelar roll.

Betrakta ett \textbf{break}: ledningen hålls låg långt längre än en ram, vilket är hur en sändare
som tappar strömmen ser ut. Ett nivåtest återaktiveras i samma ögonblick som den föregående ramens
stoppbit bedöms, så mottagaren tillbringar hela breaket med att marschera genom ram efter ram av
idel nollor. Var och en slutar på en låg stoppbit och förkastas korrekt, och den delen är i sin
ordning. Problemet är den ram som är i luften när breaket \emph{slutar}: dess stoppbit samplas efter
att ledningen gått tillbaka hög, så det är en fullkomligt välformad ram så vitt mottagaren kan
avgöra, och en byte gjord av ingenting lämnas över som riktiga data. Flanktestet återaktiveras inte
alls under breaket, eftersom det efter den första fallande flanken inte kommer någon till förrän
ledningen har hämtat sig.

Det är hela skillnaden, den kostar en vippa, och det är därför en mottagare som ”fungerar” på rena
data ändå kan lämna över skräp första gången en kabel dras ur.

\task{Skiftet}
Åtta databitar kommer med den mest signifikanta först, så SIPO-idiomet från \secref{c6:sec:sipo}
klarar det helt utan indexaritmetik:

\begin{vhdlcode}
frame <= frame(6 downto 0) & rx_s2;
\end{vhdlcode}

\code{bit_count} räknar då bara \emph{hur många} bitar som kommit, aldrig \emph{var} de hamnar. Att
i stället lagra med \code{frame(bit_count) <= rx_s2} lägger den först ankomna biten i bit 0, vilket
vänder byten: det är den ordning med minst signifikant först som en riktig UART använder, och det
skulle vara rätt val om den här mottagaren måste prata med en sådan. Här ligger MSB först på
ledningen, så skiftet är det som matchar.

Fyra saker avgör om det här fungerar:
\begin{itemize}
  \item \textbf{Allt är villkorat på \code{baud_tick}.} Ett tillstånd som agerar på en vanlig
    klockflank körs sexton gånger per tick och ungefär femtusen gånger per bit.
  \item \textbf{Varje tillstånd behöver sin ”inte än”-väg.} Om ticket du väntar på inte har kommit,
    eller räkningen inte har nått sitt mål, måste maskinen \textbf{stå kvar där den är}. Att skicka
    den någon annanstans i \code{else} är det överlägset vanligaste sättet den här designen fallerar
    på, och den fallerar fullständigt snarare än subtilt: maskinen hamnar i att studsa fram och
    tillbaka mellan två tillstånd och når aldrig det tredje.
  \item \textbf{Se upp med ettfelet i räknaren.} Med \code{ticks} som börjar på \code{0} och ökar
    vid varje tick är det $n$:te ticket \code{ticks = n - 1} i det ögonblick du testar det. Om du
    jämför före eller efter ökningen avgör vilket tick du faktiskt samplar på, och att vara ett
    eller två tick sen går att överleva medan åtta inte gör det.
  \item \textbf{Timern frilöper, så din fas kommer ur \code{ticks}, aldrig ur timern.} Startflanken
    kommer när sändaren skickar den, vilket är någonstans inuti en tickperiod mottagaren inte valt,
    så det första ticket efter den flanken kan komma allt från ett helt tick till nästan ingen tid
    senare. Att nollställa \code{ticks} på vägen ut ur vilotillståndet är det som fasar om
    samplingen mot startbiten. Den överblivna bråkdelen av ett tick är priset för att aldrig stoppa
    timern, och det är billigt: en sextondel av en bit, mot den halva bit marginal del \textbf{b)}
    ber dig räkna ut.
\end{itemize}

\xpart{d} Kör testbänken. Den behöver fyra designfiler, delblocken först, och inget av de tre
delblocken är utdelat i övningskatalogen: alla tre är moduler du skrivit.

\begin{shell}
cd lectures/L08/exercises/uart_rx8
cp ../../../L04/exercises/reset_sync/reset_sync.vhd .     # the three you wrote yourself
cp ../../../L05/exercises/sync/sync.vhd .
cp ../../../L07/exercises/timer/timer.vhd .
ghdl -a --std=93 sync.vhd reset_sync.vhd timer.vhd uart_rx8.vhd uart_rx8_tb.vhd
ghdl -e --std=93 uart_rx8_tb
ghdl -r --std=93 uart_rx8_tb --assert-level=error --stop-time=10ms
\end{shell}

Den agerar sändare, och kontrollerar:
\begin{itemize}
  \item två rena ramar, som kommer fram som rätt byte, utan något \code{frame_err}.
  \item en vilande ledning, före, mellan och efter dem, som inte producerar något.
  \item en ram vars stoppbit är \textbf{låg}: \code{frame_err} måste pulsa och ingen byte får lämnas
    över.
  \item ledningen hållen låg långt längre än en ram, vilket är ett \textbf{break}: varje ram
    mottagaren tror sig se där slutar med en låg stoppbit, så ingen av dem är en byte.
  \item en ram där varje databit bär sitt värde bara i \textbf{mitten} av bitperioden och motsatt
    nivå på var sida. En mottagare som samplar nära tick 8 läser den perfekt; en som samplar nära en
    gräns läser grannbiten.
  \item att \code{byte_valid} och \code{frame_err} var för sig aldrig är höga två cykler i rad.
\end{itemize}

\note[Tips]Bygg den i etapper och kör testbänken efter varje, i stället för att skriva alla fyra
tillstånden och felsöka dem på en gång. Få först maskinen att lämna vilotillståndet på en startbit
och nå \code{STATE_DATA}; kommer den aldrig dit är det ”inte än”-vägen du ska titta på. Sampla sedan
en databit och kontrollera att det är rätt bit innan du bryr dig om de andra sju. Sedan stoppbiten
och de två pulserna.

\xpart{e} Ha nu sönder den med flit, och fundera på vad resultatet bevisar och inte bevisar:
\begin{itemize}
  \item Ta bort stoppbitstillståndet, och lämna över byten så snart den åttonde databiten är inne.
    Förutsäg vad som borde gå fel, och kör den sedan. Två kontroller fångar det här, och det är värt
    att veta vilken som utlöser först: ramen med låg stoppbit höjer inte längre \code{frame_err},
    och den lämnar över en byte som skulle ha kastats.
  \item Nu andra halvan av lärdomen. Att passera är den svagare signalen, så hitta en sabotering
    testbänken \textbf{inte} fångar. Här är en: \textbf{ta bort omkontrollen mitt i startbiten}, så
    att \code{STATE_START} går vidare till \code{STATE_DATA} vid tick 8 utan att bekräfta att
    ledningen fortfarande är låg. Testbänken lägger aldrig någon glitch på en vilande ledning, så
    ingenting märker det, och på riktig hårdvara startar en enda brusspik en spökram. Läs
    \file{uart_rx8_tb.vhd} och säg vad du skulle lägga till för att fånga det.
  \item Här är en andra: \textbf{sampla vid tick 3 i stället för tick 8.} Räkna ut varför testbänken
    fortfarande passerar, och vad du har gett bort. Svaret ligger i marginalberäkningen i del
    \textbf{b)}.
  \item Och en tredje, som är den viktigaste av de tre: \textbf{ta bort synkroniseraren för
    \code{rx}} och läs pinnen direkt. Varje kontroll passerar fortfarande, och det kommer den alltid
    att göra, oavsett vad som läggs till i testbänken. Säg varför, i termer av vad en simulator
    modellerar och vad den inte modellerar, och läs sedan om \secref{c4:sec:mtbf}. Det här är den
    enda sortens bugg i boken som testning inte kan nå, vilket är precis varför \chapref{c4:ch} ger
    dig en regel att följa snarare än ett symptom att leta efter.
  \item Med några meningar: vad lovar egentligen ”passerar sin testbänk”, och vad lovar det inte?
\end{itemize}

\modulefig[0.74]{lectures/L08/appendix/images/uart_rx8.png}

\begin{selfcheck}
\textbf{Självkontroll.} Döp din entitet till \code{uart_rx8}, med generic
\code{OVERSAMPLE_TICK_COUNT} (\code{natural}), ingångarna \code{clock}, \code{reset_n}, \code{rx},
och utgångarna \code{data_out} (\code{std_logic_vector(7 downto 0)}), \code{byte_valid} och
\code{frame_err}, deklarerade i den ordningen. Dess testbänk finns i \file{exercises/uart_rx8/}, och
inget annat: \textbf{kopiera in dina egna \code{reset_sync}, \code{sync} och \code{timer}}, som del
\textbf{d)} visar. En projektkatalog innehåller allt den byggs av, och att sätta ihop den är en del
av övningen.
\end{selfcheck}

\exercise{\code{uart_tx8}: sändaren}{VHDL}
\label{c8:ex:tx}
Den andra halvan av länken: skriv \code{uart_tx8}, en sändare som skickar en byte i samma format som
Övning~\ref{c8:ex:rx} tar emot.

En sändare är den enklare av de två, och skälet är värt att förstå innan du börjar:
\begin{itemize}
  \item Mottagaren måste \emph{återskapa} bittajmingen ur en startflank den inte styrde över, vilket
    är varför den samplade mitt i biten och varför dess timer gick på en sextondels bitperiod.
  \item Sändaren \textbf{definierar} tajmingen. Ingenting behöver återskapas, så timern går på en
    hel bitperiod och maskinen håller helt enkelt varje bit på ledningen under en timeout.
\end{itemize}

Den sänder den mest signifikanta biten först, vilket matchar mottagaren i Övning~\ref{c8:ex:rx}
snarare än en riktig UART, av skälet som anges där.

\task{Entiteten}

\begin{booktable}{@{}p{9em}p{4.4em}p{3.4em}L@{}}
  \thead{Generic} & \thead{Typ} & \thead{Förval} & \thead{Beskrivning} \\
  \midrule
  \code{BIT_TICK_COUNT} & \code{natural} & \code{5207} & Tickantalet för \code{timer} för en
    \textbf{hel} bitperiod, 9600 baud vid \code{50 MHz}. Sexton av Övning~\ref{c8:ex:rx}:s
    översamplingstick, så när som på den avrundning dess del \textbf{b)} ber dig räkna ut, så att en
    \code{uart_tx8} och en \code{uart_rx8} byggda med de matchande förvalen pratar i samma takt. \\
\end{booktable}

\begin{booktable}{@{}p{5.4em}p{3.4em}p{10.8em}L@{}}
  \thead{Port} & \thead{Riktn.} & \thead{Typ} & \thead{Beskrivning} \\
  \midrule
  \code{clock} & in & \code{std_logic} & Systemklocka. \\
  \code{reset_n} & in & \code{std_logic} & Asynkron, aktiv låg; genom \code{reset_sync}, som
    alltid. \\
  \code{data_in} & in & \code{std_logic_vector(7 downto 0)} & Byten som ska skickas. \\
  \code{send} & in & \code{std_logic} & En begäran om att sända \code{data_in}. Ignoreras medan
    sändaren är upptagen. \\
  \code{tx} & out & \code{std_logic} & Den seriella ledningen. \textbf{Vilar hög}, även under
    reset. \\
  \code{busy} & out & \code{std_logic} & Hög från ögonblicket en ram startar tills stoppbiten är
    över. \\
\end{booktable}

\xpart{a} Rita tillståndsmaskinen först. Fyra tillstånd är den naturliga uppdelningen: vila,
startbit, databitar, stoppbit, och vart och ett håller ledningen på en definierad nivå under en
bitperiod.
\begin{itemize}
  \item Markera vad som driver \code{tx} i varje tillstånd. Lägg märke till att \code{tx} aldrig
    beror på \code{send} eller någon annan primär ingång: det är tillståndet, plus skiftregistrets
    aktuella utgång. Det gör det här till en Mooremaskin, som allt annat i den här boken.
  \item Markera var \code{busy} kommer ifrån. Det är en enda jämförelse mot tillståndet.
  \item Bestäm var byten fångas. Den måste låsas när ramen startar, inte läsas kontinuerligt ur
    \code{data_in}: producenten står fritt att ändra \code{data_in} i samma ögonblick som
    \code{busy} går hög, och en verklig sådan kommer att göra det.
\end{itemize}

\xpart{b} Skifta ut byten med den mest signifikanta biten först, i maskinen själv, spegelbilden av
vad mottagaren gör: lås in \code{data_in} i en signal \code{frame} när ramen startar, driv \code{tx}
från den översta biten i \code{frame} medan du är i datatillståndet, och skifta \code{frame} ett steg
uppåt per bitperiod, med räkning på de åtta bitar som går ut.

Det här skrev du som en modul i Övning~\ref{c6:ex:shiftregs}, och \code{piso8} är värd att läsa om
för idiomets skull. Att instansiera den här skulle fungera, men maskinen äger redan den biträknare
ett PISO-register behöver, så skiftet blir en enda rad inuti det tillstånd där det hör hemma. Så är
båda halvorna av den här länken byggda, och samma avvägning återkommer i \chapref{c14:ch}, där
\code{tx_shift_reg} är en egen modul just för att \code{can_controller} \emph{inte} redan äger den
räknare den behöver.

\xpart{c} Implementera den, med \code{reset_sync} instansierad för reset och \code{timer} med
\code{BIT_TICK_COUNT} skickat rakt igenom.

Här \textbf{är} timerns \code{enable} driven av tillståndet, hög bara medan en ram är i luften, och
det är det enda stället där den här designen medvetet skiljer sig från Övning~\ref{c8:ex:rx}.
Mottagaren måste låta sin timer frilöpa eftersom den inte styr när en ram börjar: den kan bara fasa
om sin egen tickräkning när startflanken redan har kommit. Sändaren har det motsatta problemet. Den
\emph{bestämmer} när ramen börjar, så den kan starta timern i det ögonblicket och få en första
bitperiod som är exakt lika lång som alla andra. Låt den frilöpa här i stället, så blir startbiten
kort med precis så mycket av en tickperiod som råkade återstå när \code{send} kom, vilket är en
defekt i det enda en sändare är ansvarig för.

Två krav testbänken är sträng med:
\begin{itemize}
  \item \textbf{\code{send} ignoreras medan sändaren är upptagen.} Att ladda om mitt i en ram skulle
    starta om byten och korrumpera det som redan ligger på ledningen. En begäran som kommer under en
    ram kastas, den köas inte.
  \item \textbf{\code{busy} ligger kvar hög genom hela stoppbiten}, inte bara databitarna. Ramen är
    inte över förrän ledningen hållits hög under den sista bitperioden.
\end{itemize}

\xpart{d} Kör testbänken. Den agerar mottagare: den väntar på din startbit, samplar mitten av varje
bitperiod, och kontrollerar ramningen, byten, handskakningen med \code{busy}, och att ett
\code{send} som matas in två bitperioder in i en ram ignoreras.

\begin{shell}
cd lectures/L08/exercises/uart_tx8
cp ../../../L04/exercises/reset_sync/reset_sync.vhd .     # the two you wrote yourself
cp ../../../L07/exercises/timer/timer.vhd .
ghdl -a --std=93 reset_sync.vhd timer.vhd uart_tx8.vhd uart_tx8_tb.vhd
ghdl -e --std=93 uart_tx8_tb
ghdl -r --std=93 uart_tx8_tb --assert-level=error --stop-time=10ms
\end{shell}

\xpart{e} Svara, i löpande text:
\begin{itemize}
  \item En sändare som hoppar över stoppbiten helt lämnar ändå ledningen hög efteråt, eftersom vila
    och stopp är samma nivå. Vad går egentligen fel, och när visar det sig först? Testbänken fångar
    den här: vilken av dess kontroller gör det, och varför hade det aldrig gått att fånga genom att
    bara titta på \code{tx}?
  \item Din mottagare i Övning~\ref{c8:ex:rx} samplar mitt i varje bit; din sändare ändrar ledningen
    vid flankerna. Om båda gick på kristaller som skilde sig med 1~\%, ungefär hur många bitar in i
    en ram skulle mottagarens samplingspunkt driva in i fel bit? Vad säger det dig om varför
    seriella format har en startbit per byte, snarare än en enda i början av ett långt meddelande?
\end{itemize}

\modulefig[0.66]{lectures/L08/appendix/images/uart_tx8.png}

\begin{selfcheck}
\textbf{Självkontroll.} Döp din entitet till \code{uart_tx8}, med generic \code{BIT_TICK_COUNT}
(\code{natural}), ingångarna \code{clock}, \code{reset_n}, \code{data_in}
(\code{std_logic_vector(7 downto 0)}), \code{send}, och utgångarna \code{tx}, \code{busy},
deklarerade i den ordningen. Dess testbänk finns i \file{exercises/uart_tx8/}, och inget annat:
kopiera in dina egna \code{reset_sync} och \code{timer}, som del \textbf{d)} visar.
\end{selfcheck}
